% gz-p7-03-1: 散点图 + 线性回归——15 个散点，正相关，叠加回归直线
\documentclass[tikz,border=4pt]{standalone}
\usepackage{ctex}
\usepackage{amsmath}
\usepackage{pgfplots}
\pgfplotsset{compat=1.18}

\begin{document}
\begin{tikzpicture}
\begin{axis}[
  width=10cm, height=7cm,
  axis lines=center,
  xlabel={$x$}, ylabel={$y$},
  every axis x label/.style={at={(axis cs:10.5,0)}, anchor=west},
  every axis y label/.style={at={(axis cs:0,10.5)}, anchor=south},
  xmin=0, xmax=10.5,
  ymin=0, ymax=10.5,
  xtick={2,4,6,8,10},
  ytick={2,4,6,8,10},
  tick style={color=black},
  ticklabel style={font=\small},
  clip=true,
]
  % 15 个散点（正相关）
  \addplot[only marks, mark=*, mark size=2.0pt, blue] coordinates {
    (1, 1.6) (1.5, 2.4) (2, 1.9) (2.5, 3.2) (3, 3.0)
    (3.5, 3.8) (4, 4.6) (4.5, 4.3) (5, 5.4) (5.5, 5.7)
    (6, 5.9) (6.5, 7.2) (7, 6.8) (7.5, 7.7) (8, 8.3)
  };

  % 回归直线 y = 0.95 x + 0.3
  \addplot[red, thick, domain=0.5:8.5] {0.95*x + 0.3};
  \node[right, red] at (axis cs:7.5, 7.4) {$\hat{y}=\hat{b}x+\hat{a}$};

  % 注释
  \node[font=\footnotesize, blue, align=left] at (axis cs:2, 9)
    {散点呈正相关\\拟合直线斜率 $>0$};
\end{axis}
\end{tikzpicture}
\end{document}
